An observation too tentative for an explanation can already warrant a check. When people and models trade roles in spotting and testing leads, training for safety\cite{SAFE,BIAN} could alter which questions they pursue. Here we show that swapping only the fitted Qwen3-8B writer on identical records altered whether a note justified the next decision. Request-level protection reduced pre-filter disclosures in a separate office test. At 650 tokens it yielded 8/24 usable notes versus 18/24 under ordinary lawful help on eight new public tasks across three paired seeds. On five archived histories preselected by a masked assessor for sufficient grounds, the counts reversed to 4/5 versus 0/5 (4/24 versus 0/24 overall). In one case the ordinary-help investigator commissioned a scan but omitted its counterexample at scale 9. The protective writer conveyed it from the same record. One recipient showed this check first. Another proved the same correction from the public problem alone. All 18 frontier-model runs on six selected problems reached justified decisions. In all 72 assignments, ordinary workflows reproduced the corresponding field forecast (eight fields, nine configurations each). Whether the same training that aids current work leaves needed questions or methods unformed or unrecovered in time at comparable cost remains open.
